\documentclass[12pt]{report}
% Template Tesi in Informatica, ispirata al template "Tesi di laurea - Universit� di Pisa" di Simone Schirinzi (https://www.overleaf.com/latex/templates/tesi-di-laurea-universita-di-pisa/rwdcqtqwftpg). 

% Carattere dimensione 12

% Per la stampa fronte-retro sostituire con:
% \documentclass[12pt, twoside]{report}

% Margini (4cm a sx, 2.5cm a dx, 2.5cm in alto, 2.5cm in basso)
\usepackage[top=2.5cm, bottom=2.5cm, left=4cm, right=2.5cm, head=21.75pt, centering]{geometry}

% Per la stampa fronte-retro sostituire con: 
% \usepackage[top=2.5cm, bottom=2.5cm, inner=4cm, outer=4cm, right=2.5cm, centering]{geometry}

% Interlinea
\linespread{1.5}

% Librerie utili
\usepackage[italian]{babel} % applicazione regole di scrittura per la lingua italiana 
\usepackage[utf8]{inputenc} % codifica UTF-8
\usepackage{scrlayer-scrpage} % stili pagina per il frontespizio
\usepackage[hidelinks]{hyperref} % collegamenti ipertestuali
\ifoot[]{}
\cfoot[]{}
\ofoot[\pagemark]{\pagemark}
\pagestyle{scrplain}
% \usepackage{mathptmx} % font Times New Roman (simile)
\usepackage{graphicx} % inserimento di immagini
\usepackage{csquotes} % per le citazioni "in blocco"
\usepackage{svg}
% \usepackage[backend=biber, sorting=nty, ]{biblatex} % bibliografia con pacchetto biblatex (https://ctan.org/pkg/biblatex?lang=en)
\appto{\bibsetup}{\raggedright}

\newcommand{\quotes}[1]{``#1''}

\usepackage{titlesec} % per la formattazione dei titoli delle sezioni, capitoli etc.
\usepackage{float} % per il posizionamento delle immagini

\usepackage{listings} % per il codice di programmazione
% Fonte https://en.wikibooks.org/wiki/LaTeX/Source_Code_Listings. Per la lista di sintassi riconosciute.
\renewcommand{\lstlistingname}{Code}% Listing -> Codice
\usepackage{xcolor}  % stile del codice
\definecolor{mygreen}{rgb}{0,0.6,0}
\definecolor{mygray}{rgb}{0.5,0.5,0.5}
\definecolor{mymauve}{rgb}{0.58,0,0.82}
\definecolor{darkgray}{rgb}{.4,.4,.4}
\definecolor{navy}{HTML}{000080}
\definecolor{purple}{rgb}{0.65, 0.12, 0.82}
\definecolor{codepurple}{rgb}{0.58,0,0.82}
\definecolor{backcolour}{rgb}{0.95,0.95,0.92}

% Stili configurabili del codice (lslisting) 
\lstset{ %
  belowcaptionskip=0.5em,
  backgroundcolor=\color{backcolour}, % choose the background color; you must add \usepackage{color} or \usepackage{xcolor}
  basicstyle=\footnotesize, % the size of the fonts that are used for the code
  breakatwhitespace=false, % sets if automatic breaks should only happen at whitespace
  breaklines=true, % sets automatic line breaking
  captionpos=b, % sets the caption-position to bottom
  commentstyle=\color{mygreen}, % comment style
  deletekeywords={...}, % if you want to delete keywords from the given language
  escapeinside={\%*}{*)}, % if you want to add LaTeX within your code
  extendedchars=true, % lets you use non-ASCII characters; for 8-bits encodings only, does not work with UTF-8
  frame=single, % adds a frame around the code
  keepspaces=true, % keeps spaces in text, useful for keeping indentation of code (possibly needs columns=flexible)
  keywordstyle=\color{codepurple}, % keyword style
  language=Python, % the language of the code
  morekeywords={*,...}, % if you want to add more keywords to the set
  numbers=left, % where to put the line-numbers; possible values are (none, left, right)
  numbersep=5pt, % how far the line-numbers are from the code
  numberstyle=\tiny\color{mygray}, % the style that is used for the line-numbers
  rulecolor=\color{black}, % if not set, the frame-color may be changed on line-breaks within not-black text (e.g. comments (green here))
  showspaces=false, % show spaces everywhere adding particular underscores; it overrides 'showstringspaces'
  showstringspaces=false, % underline spaces within strings only
  showtabs=false, % show tabs within strings adding particular underscores
  stepnumber=1, % the step between two line-numbers. If it's 1, each line will be numbered
  stringstyle=\color{mymauve}, % string literal style
  tabsize=2, % sets default tabsize to 2 spaces
  title=\lstname % show the filename of files included with \lstinputlisting; also try caption instead of title
}

% END of listing package 

% Esempio riconoscimento sintassi JavaScript 

\lstdefinelanguage{JavaScript}{
  keywords={typeof, new, true, false, catch, function, return, null, catch, switch, var, const, let, if, in, while, do, else, case, break},
  keywordstyle=\color{blue}\bfseries,
  ndkeywords={class, export, boolean, throw, implements, import, this},
  ndkeywordstyle=\color{darkgray}\bfseries,
  identifierstyle=\color{black},
  sensitive=false,
  comment=[l]{//},
  morecomment=[s]{/*}{*/},
  commentstyle=\color{mygreen}\ttfamily,
  stringstyle=\color{red}\ttfamily,
  morestring=[b]',
  morestring=[b]"
}

\lstset{
   language=JavaScript,
   extendedchars=true,
   basicstyle=\footnotesize\ttfamily,
   showstringspaces=false,
   showspaces=false,
   numbers=left,
   numberstyle=\footnotesize,
   numbersep=9pt,
   tabsize=2,
   breaklines=true,
   showtabs=false,
   captionpos=b
}

% Formato delle intestazioni
\titleformat{\chapter}[block]
  {\normalfont\LARGE\bfseries}{\thechapter.}{0.5em}{\LARGE}
\titlespacing*{\chapter}{0pt}{-20pt}{25pt}

\begin{document}

% Frontespizio
\begin{titlepage}
  \begin{figure}
    \centering\includegraphics[scale=0.5]{immagini/cherubino_pant541.png}
  \end{figure}

  \begin{center}
    {\LARGE{ Corso di Laurea in Informatica \\ }}
    \vspace{2cm}
    {\Large { TESI DI LAUREA }}\\
    \vspace{2cm}
    {\Large { Algoritmi genetici per Explainable AI }}
  \end{center}

  \vspace{2cm}

  \begin{minipage}[t]{0.47\textwidth}
    {\large{Relatori:\\ Riccardo Guidotti}}
    \vspace{0.5cm}
    {\large{\\Correlatore:\\ Marco Danelutto}}
  \end{minipage}\hfill\begin{minipage}[t]{0.47\textwidth}\raggedleft
    {\large{Candidato: \\ Federico Bustaffa\\ }}
  \end{minipage}

  \vspace{25mm}

  \centering{\large{\bf ANNO ACCADEMICO 2024/2025 }}
\end{titlepage}
% Fine frontespizio

\tableofcontents
\thispagestyle{empty}

\listoffigures

\thispagestyle{empty}
\clearpage
\setcounter{page}{1}
\addtocontents{toc}{\protect\thispagestyle{empty}}
\addcontentsline{toc}{chapter}{Introduzione} % Capitolo non numerato

\chapter*{Introduzione}

Il progetto di tesi verte sull'implementazione di una libreria di algoritmi
genetici, in grado di sfruttare architetture multi-core per il calcolo
parallelo. La necessità di un algoritmo genetico con questa struttura nasce in
un contesto di explainable AI, in particolare si fa riferimento al progetto
\textit{LORE}~\cite{guidotti2018LORE}, il quale utilizza un'implementazione
fornita dalla libreria di algoritmi genetici \textit{DEAP}~\cite{fortin2012DEAP}.

\textit{LORE} si propone di generare spiegazioni a decisioni o predizioni fatte
da modelli di machine learning, spesso difficili da interpretare per via della
loro elevata complessità, come ad esempio reti neurali profonde o foreste
casuali. In una prima fase, del metodo si sfrutta un algoritmo genetico per la
generazione di dati sintetici con determinate caratteristiche, fondamentali in
seguito per la costruzione delle spiegazioni.

Tale approccio risulta però essere particolarmente dispendioso, soprattutto in
casi in cui la quantità di dati sintetici che si desidera generare è molto
grande oppure quando il modello è particolarmente lento in fase di predizione
(o entrambe).

Da qui la necessità di implementare una libreria di algoritmi genetici in grado
di produrre risultati molto simili a \textit{DEAP} ma più performante e in
grado di sfruttare al meglio architetture multi-core.

Dato che \textit{LORE} e \textit{DEAP}, così come la maggior parte delle
librerie di machine learning, sono implementate (o forniscono un'API) in Python,
la prima fase del lavoro si è focalizzata sull'esplorazione e la ricerca di
possibili framework, in grado di lavorare in sinergia con Python. Le opzioni
esplorate vanno da librerie standard di Python e moduli di terze parti, passando
per l'implementazione di estensioni in linguaggi a basso livello come C/C++,
fino ad arrivare alla sperimentazione di feature sperimentali come
\textit{subinterpreters} e una delle prime versioni \textit{free-threaded}
proposta con Python 3.13.

Una volta implementata la libreria si è passati ad una fase di test per
valutarne la correttezza, impiegando l'algoritmo nella risoluzione di problemi
ben noti in letteratura facilmente rappresentabili tramite grafici. Tra questi
abbiamo il problema dello zaino, del commesso viaggiatore e un semplice caso di
regressione lineare. Si è infine testata la correttezza sul caso d'uso di
\textit{LORE} precedentemente descritto.

Il lavoro si è concluso con due fasi di test, la prima riguardante la qualità
delle soluzioni prodotte e la seconda riguardante le performance. Come
anticipato, l'obiettivo della tesi è stato quello di produrre qualcosa che
funzionasse qualitativamente come \textit{DEAP} ma che fosse più performante.

La libreria implementata sembra soddisfare le aspettative fornendo risultati
qualitativi in linea con \textit{DEAP} ma dimostrando prestazioni e capacità
di sfruttare architetture multi-core generalmente migliori.


\clearpage

\chapter{Contesto}

Come anticipato nell'introduzione, il contesto in cui si va a collocare il
lavoro di tesi è molteplice. Abbiamo infatti a che fare con algoritmi genetici,
calcolo parallelo ed accenni di explainable AI.

\section{Algoritmi genetici}

Gli algoritmi genetici fanno parte delle possibili tecniche euristiche adottate
nell'ambito dell'ottimizzazione. Costituiscono infatti una valida scelta in
molti casi in cui non si conosce un algoritmo di ottimizzazione efficiente per
un problema specifico e non si è interessati alla ricerca di soluzioni ottime
per esso.

Il loro comportamento si ispira ad un modello evolutivo in cui gli individui
migliori di una popolazione sopravvivono e si riproducono, portando quindi la
popolazione iniziale ad uno stadio evolutivo migliore con il passare delle
generazioni.

Le possibili implementazioni sono molto varie ma una delle più comuni prevede
la rappresentazione di ogni individuo tramite un cromosoma. Questo è
generalmente strutturato come un vettore o una stringa, i cui elementi
rappresentano una possibile soluzione al problema o una sua codifica. Una delle
implementazioni più semplici è composta dalle seguenti fasi:
\begin{enumerate}
    \item Generazione: viene generata in modo casuale una popolazione i cui
          individui rappresentano possibili soluzioni al problema.
    \item Selezione: vengono selezionati $k$ individui utilizzando metodi
          basati sul valore di fitness.
    \item Crossover: i cromosomi di due o più individui vengono combinati per
          generare nuovi individui che condividono parte del materiale genetico
          dei genitori.
    \item Mutazione: vengono applicate piccole variazioni casuali al cromosoma
          di un individuo, permettendo così l'esplorazione di nuove soluzioni e
          favorendo così la diversità genetica all'interno della popolazione.
    \item Valutazione: ad ogni individuo viene attribuito un valore di
          \textit{fitness} in base alla qualità della soluzione che questo
          rappresenta.
    \item Rimpiazzo: viene attuata una politica di rimpiazzo (in alcuni casi
          chiamata \quotes{di sopravvivenza}), che prevede la scelta degli
          individui da mantenere da una generazione all'altra.
\end{enumerate}
Tutte le fasi tranne la prima sono eseguite in un ciclo che termina quando
viene soddisfatto un criterio di convergenza, che può essere definito in
termini dei valori di fitness della popolazione, della sua biodiversità o, più
semplicemente, definendo un limite massimo al numero delle iterazioni.
L'algoritmo risulta quindi essere strutturato come segue
\begin{figure}
    \centering
    \includesvg[scale=0.7]{immagini/genetic_algorithm.svg}
    \caption{Struttura di un algoritmo genetico semplice}
    \label{fig:simple_ga}
\end{figure}
Vista la natura stocastica dell'algoritmo, è inoltre buona norma avere a
disposizione una struttura dati in grado di memorizzare le migliori $k$
soluzioni mai registrate tra le varie generazioni.

\subsection{Calcolo parallelo}

L'obiettivo principale del lavoro svolto è stato quello di rimodellare la
struttura appena descritta, di modo da riuscire ad eseguire in parallelo le
fasi di crossover, mutazione e valutazione, le quali rappresentavano
(soprattutto la valutazione) il vero collo di bottiglia per le prestazioni nel
caso d'uso di riferimento. Il modello di calcolo ha quindi assunto la seguente
forma
\begin{figure}
    \centering
    \includesvg[scale=0.65]{immagini/parallel_ga.svg}
    \caption{Algoritmo genetico parallelo}
    \label{fig: parallel_ga}
\end{figure}
che mira ad implementare un paradigma di tipo \textit{map}. Queste fasi
si prestano bene a questo tipo di paradigma poiché consistono
nell'applicare una funzione in modo indipendente ad ogni elemento di una
lista. In realtà crossover e mutazione vengono applicate con una certa
probabilità, ma una volta terminate, i nuovi individui vengono
sicuramente tutti valutati.

\subsubsection{Calcolo parallelo in Python}

Implementare modelli di calcolo parallelo in Python può non essere così
immediato come in altri linguaggi come C, C++ o Java. In Python non è infatti
possibile sfruttare il multithreading per compiti computazionalmente pesanti,
a causa del \textit{GIL} (Global Interpreter Lock). Come suggerisce il nome,
si tratta di una variabile di mutua esclusione che ogni thread deve acquisire
per poter eseguire il proprio codice. Cercare quindi di sfruttare il
multithreading per questo tipo di compiti non porterà alcun beneficio poiché
i thread verranno eseguiti uno per volta. Potrebbe anzi verificarsi un
deterioramento delle prestazioni dovuto all'overhead di scheduling e
sincronizzazione che il multithreading comporta.

Al momento ciò che viene considerato lo standard per risolvere questo tipo di
problema è l'impiego di un paradigma multi-processo. In questo modo si hanno
più interpreti che operano indipendentemente l'uno dall'altro e su spazi di
memoria differenti. Questo approccio richiede però maggiori risorse sia in
termini di memoria, sia in termini di tempo poiché la condivisione dati avviene
tramite meccanismi di streaming di byte (come pipe o socket) e non in memoria
condivisa come per il multithreading.

Esistono però altri metodi che, sotto determinate ipotesi, permettono di
sfruttare il parallelismo senza dover ricorrere all'utilizzo di processi
multipli. Con le nuove versioni di Python sono state inoltre proposte nuove
funzionalità (ancora sperimentali) che mirano proprio a risolvere o mitigare il
problema. Parte del lavoro è stato infatti volto all'esplorazione di tali
opzioni per garantire la massima efficienza possibile, garantendo però un certo
livello di espressività.

\section{Explainable AI}

L'explainable AI, abbreviato XAI, è una branca della ricerca volta a fornire
metodi per interpretare e \textit{spiegare} predizioni o decisioni fatte da
determinati modelli di machine learning, chiamati \textit{black-box}. Di questi
si conoscono architettura e modello matematico, ma allo stato dell'arte possono
arrivare ad avere una quantità di parametri talmente elevata da rendere
impossibile per un essere umano capire come e quanto i singoli parametri
influiscano nelle decisioni prese dal modello.

Al fine di riuscire ad interpretare le scelte fatte dai modelli
\textit{black-box}, sono stati sviluppati vari metodi. Alcuni di questi cercano
di comprendere la struttura globale del modello predittivo, per esempio cercando
di capire la forma dei confini di classificazione. Altri metodi si concentrano
invece sul fornire spiegazioni alle singole decisioni prese dal modello su
specifiche istanze dei dati, fornendo quindi spiegazioni locali. I metodi di
XAI possono inoltre fornire \textit{controfattuali}, ossia spiegazioni su come
modificare i dati i dati in input per cambiare le decisioni prese dal modello.

% approfondimento possibili approcci e metodi...


\clearpage

\chapter{Stato dell'arte}

% 1. Iniziare a definire qui le possibili opzioni per l'implementazione di
% algoritmi paralleli in Python? Come multiprocessing, Cython, C/C++ extensions,
% e subinterpreters

% 2. PARALLELISMO: Spiegare solo il modello di parallelismo offerto dalla
% combinazione di DEAP e multiprocessing o anche di parallelismo generale (map)?

Allo stato dell'arte, un'implementazione di quanto citato nel capitolo
precedente è già presente. Abbiamo infatti \textit{LORE}~\cite{guidotti2018LORE}
come metodo di XIA, un'implementazione di algoritmo genetico fornita da
\textit{DEAP}~\cite{fortin2012DEAP}, e un metodo per parallelizzare tale
algoritmo, tramite la combinazione di \textit{DEAP} e del modulo
\verb|multiprocessing| di Python.

\section{LORE}

Il metodo proposto da \textit{LORE}~\cite{guidotti2018LORE} rientra nella
categoria dei metodi locali. Opera infatti sulle singole istanze dei dati,
lavorando sotto l'assunzione che i confini di classificazione locali siano
generalmente meno complessi di quelli globali, riuscendo così ad ottenere % verificare
spiegazioni più accurate.

Il modo in cui opera \textit{LORE} prevede l'esecuzione di un algoritmo genetico
su ognuna delle istanze classificate dal modello e di cui si vogliono ottenere
fattuali e controfattuali. In una prima fase si vuole quindi generare un
vicinato sintetico i cui individui siano classificati allo stesso modo
dell'istanza di riferimento, si cerca poi di generare individui sintetici
classificati diversamente da essa, minimizzando in entrambi i casi la distanza
tra i vicini sintetici e l'istanza stessa.

Una volta generati, i dati sintetici vengono usati per allenare un albero di
decisione contenente delle regole che andranno poi a costituire la spiegazione
della scelta fatta e le possibili modifiche sui dati, di modo da cambiare tale
scelta. Le diramazioni di un albero decisionale corrispondono infatti a valori
soglia delle feature che determinano una scelta piuttosto che un'altra. Seguendo
quindi i cammini che portano alla stessa decisione effettuata dal modello si
ottiene la spiegazione di tale scelta, seguendo invece cammini alternativi è
possibile capire come cambiare tale scelta.

\section{DEAP}

La libreria \textit{DEAP} fornisce un'API molto flessibile ma richiede un fase
di inizializzazione in cui vengono definiti tipi, strutture dati e funzioni
necessarie al corretto funzionamento dell'algoritmo.

\subsection{Inizializzazione}

Come prima cosa è necessario definire la struttura di un generico cromosoma, il
tipo di ciascun gene e i pesi da assegnare a ciascun parametro da ottimizzare.
A tal fine vengono messi a disposizione il modulo \verb|creator| e la classe
\verb|Toolbox|.

Tramite \verb|creator| è possibile definire nuove classi con caratteristiche
personalizzate e che si adattano alle proprie esigenze. La classe \verb|Toolbox|
serve invece per immagazzinare funzioni e relativi parametri, di modo da poter
essere utilizzati con facilità in seguito. Tramite tale classe è infatti
possibile definire operatori genetici e metodi di inizializzazione, selezione e
valutazione degli individui. Ad eccezione però della funzione di valutazione,
che è strettamente legata al problema di interesse e deve essere quindi definita
dall'utente caso per caso, \textit{DEAP} fornisce implementazioni per tutte le
componenti appena citate, lasciando quindi il solo compito di scegliere quale
si addice di più al problema in questione. Rimane comunque la possibilità di
definire ciascuna delle componenti dell'algoritmo come meglio si crede nel caso
ci siano esigenze particolari.

\subsection{Algoritmi}

Una volta definite tutte le componenti dell'algoritmo, è possibile invocare
direttamente i metodi registrati nell'istanza della classe \verb|Toolbox|,
andando a definire un comportamento personalizzato per esso. Una seconda opzione
consiste invece nel ricorrere alle implementazioni fornite da \textit{DEAP},
tramite il modulo \verb|algorithms|, che presenta diverse implementazioni
possibili, tra cui \verb|eaSimple|, che ha una struttura molto simile a quella
definita in figura \ref{fig:simple_ga} ed è l'algoritmo attualmente impiegato
da \textit{LORE}.

L'algoritmo prende in input l'istanza della classe \verb|Toolbox| precedentemente
definita, una popolazione iniziale, il numero massimo di iterazioni da compiere
e le probabilità di crossover e mutazione. \`E possibile passare anche parametri
opzionali per registrare dati statistici sui valori di fitness e una struttura
dati per registrare i migliori $k$ individui mai prodotti durante le varie
generazioni, ossia la \verb|HallOfFame|.

\subsection{Multi-processing}

La classe \verb|Toolbox| ha anche dei metodi predefiniti, tra cui il metodo
\verb|map|, la cui implementazione è data dall'omonima funzione standard di
Python. Nel caso dell'algoritmo \verb|eaSimple|, viene usata la funzione
\verb|map| registrata nell'istanza delle \verb|Toolbox| per applicare la
funzione di valutazione a ciascun individuo della popolazione.

Come per gli altri metodi è possibile registrare una qualunque funzione sotto
il nome "map" e \textit{DEAP}, al fine di ottenere l'esecuzione in parallelo,
suggerisce di registrare la funzione \verb|map| della classe \verb|Pool| del
modulo \verb|multiprocessing|.

Il modulo \verb|multiprocessing| è considerato lo standard per il calcolo
parallelo in Python quando si ha del lavoro computazionalmente intensivo da
compiere. Come suggerisce il nome sfrutta processi multipli per aggirare il
problema posto dal GIL, aumentando però l'overhead di comunicazione tra i
processi e il consumo complessivo di memoria.

Al fine di gestire i processi e la distribuzione del carico in modo più semplice
è disponibile una classe \verb|Pool|, che implementa un pool di processi in
grado di bilanciare automaticamente il carico e comunicare dati input e output
per ogni operazione. Il metodo \verb|map| della classe \verb|Pool| applica una
certa funzione a tutti gli elementi di una struttura dati iterabile passata come
argomento. Se non si specificano ulteriori argomenti, gli elementi dell'iterabile
verranno inviati uno per volta ai processi una volta che questi sono pronti. Al
contrario se si specifica il parametro \verb|chunksize| è possibile definire il
numero di elementi che un processo deve elaborare prima di poterne ricevere
altri.

\clearpage

\chapter{Metodologia}

\clearpage

\chapter{Sperimentazione}

\addcontentsline{toc}{chapter}{Conclusioni} % Capitolo non numerato
\chapter*{Conclusioni}


\bibliographystyle{plain}
\bibliography{bibliography} % Entries are in the refs.bib file
% \addbibresource{bibliography.bib}
% \printbibliography

\end{document}